\documentclass[11pt]{article}
\usepackage[a4paper,margin=18mm]{geometry}
\usepackage[no-math]{fontspec}
\setmainfont{Latin Modern Roman}
\usepackage{amsmath,amssymb,amsthm,xfrac,xspace,hyperref,graphicx,comment}
\excludecomment{DefTralics}

\providecommand{\bibliofont}{\small}
\newcommand{\ie}{i.e.,\xspace}
\newcommand{\eg}{e.g.,\xspace}
\newcommand{\etc}{etc.\xspace}
\newcommand{\cf}{cf.\xspace}
\newcommand{\resp}{resp.\xspace}
\newcommand{\smaller}{\small}
\newcommand{\Subsection}{\subsection}
\newcommand{\Subsubsection}{\subsubsection}
\newcommand{\skpt}{\smallskip}
\emergencystretch=3em
\numberwithin{equation}{section}\input{author-preamble.tex}
\begin{document}
\begin{center}{\Large Stable item 2044}\end{center}
\paragraph{Source and attribution.}
Michael Björklund, Alexander Fish, Ilya D. Shkredov, \emph{Sets of transfer times with small densities}, Journal de l’École polytechnique — Mathématiques 8 (2021), publisher version, DOI 10.5802/jep.147. \href{https://jep.centre-mersenne.org/articles/10.5802/jep.147/}{Publisher article}. Authors' text: \href{https://creativecommons.org/licenses/by/4.0/}{CC BY 4.0}.
\paragraph{Editorial problem boundary.}
Prove only Theorem 1.8 below for a totally ergodic action of a countable discrete abelian group and positive-measure sets whose lower transfer-time density equals the sum of their measures, less than one. The exact averaging-sequence assumption, compact factor correspondence and overshoot argument are retained. The return set agrees with the inverse circle difference interval modulo at most two kernel cosets, as printed. Arbitrary averaging sequences, semigroup extensions and later characterisations are outside the selected claim.
\paragraph{Difficulty (editorial): H2.}
The cited compact action-set factor theorem and Kneser critical-pair theorem do not alone identify the original action's full transfer-time set or its endpoints. The author establishes a measurable transfer-time correspondence with a quantitative overshoot inequality and uses equality of the densities to recover the actual circle factor and sets. This additional correspondence shows precisely why only the two endpoint kernel cosets can differ from the inverse circle difference interval.
\paragraph{Editorial transcription note.}
Original author language, mathematics and ordering are preserved. Publisher front matter is replaced by this independent article wrapper. Bibliography is recovered from matching cached publisher HTML, including its TeX formulas. No new proof or mathematical repair is supplied. Surrounding original results are retained for conservative context and are not extra selected corpus claims.
\clearpage
\input{author-body.tex}
\begin{thebibliography}{99}
\bibitem{BF} M. Björklund \& A. Fish - “Approximate invariance for ergodic actions of amenable groups” , Discrete Anal. (2019), article ID 6, 56 pages
\bibitem{BW} M. Boshernitzan \& M. Wierdl - “Ergodic theorems along sequences and Hardy fields” , Proc. Nat. Acad. Sci. U.S.A. 93 (1996) no. 16, p. 8205-8207
\bibitem{Bo88} J. Bourgain - “On the maximal ergodic theorem for certain subsets of the integers” , Israel J. Math. 61 (1988) no. 1, p. 39-72
\bibitem{Kneser0} M. Kneser - “Abschätzung der asymptotischen Dichte von Summenmengen” , Math. Z. 58 (1953), p. 459-484
\bibitem{Kneser} M. Kneser - “Summenmengen in lokalkompakten abelschen Gruppen” , Math. Z. 66 (1956), p. 88-110
\bibitem{La} Y. Lacroix - “Natural extensions and mixing for semi-group actions” , Publ. Inst. Rech. Math. Rennes 2 (1995), article ID 7, 10 pages
\end{thebibliography}
\end{document}
